\documentclass[11pt]{article}

% Packages
\usepackage[T1]{fontenc}
\usepackage[margin=1in]{geometry}
\usepackage{amsmath,amssymb}
\usepackage{natbib}
\usepackage{graphicx}
\usepackage{booktabs}
\usepackage{xcolor}
\usepackage{hyperref}
\usepackage{cleveref}

\title{Supplementary Material: Physical derivation of Greenland discharge model parameters}
\author{}
\date{}

\begin{document}
\maketitle

\noindent This supplement derives the key parameters of the Greenland discharge time-delay model---the ocean--discharge coupling $\gamma$, the time delay $\delta$, and the secular drift rate $r_0$---from glacier dynamics, and compares these independent estimates with the WLS-fitted values obtained by calibrating the model against the \citet{Mouginot2019} discharge record.

The discharge model (main text Eq.~\ref*{eq:greenland_dyn}) is
\begin{equation}
	\Delta H_{\mathrm{dyn}}(t) = \gamma \left[\int_{t_0}^{t} T_{\mathrm{ocean}}(t' - \delta)\, dt' - \theta\right] + r_0\, [t - \bar{t}],
	\label{eq:sup_Hdyn}
\end{equation}
where $\Delta H_{\mathrm{dyn}} = H_{\mathrm{dyn}} - \langle H_{\mathrm{dyn}} \rangle_{\mathrm{cal}}$ is the demeaned cumulative discharge, $T_{\mathrm{ocean}}(t - \delta)$ is the subsurface ocean temperature shifted by a fixed delay $\delta$, and $\theta$ and $\bar{t}$ are calibration-period means. The instantaneous discharge rate is $\dot{H}_{\mathrm{dyn}} = \gamma\, T_{\mathrm{ocean}}(t - \delta) + r_0$.

%% ================================================================
\section{Time delay $\delta$}
\label{sec:delta}

The time delay $\delta$ represents the lag between a change in subsurface ocean temperature and the resulting change in ice discharge. It combines two processes in series: (i)~warm-water renewal in glacial fjords and (ii)~the dynamic adjustment of the glacier terminus and upstream ice.

\paragraph{Fjord renewal.}
Fjord circulation replaces subsurface water through a buoyancy-driven overturning driven by subglacial discharge and fjord--shelf exchange \citep{Straneo2015}. Typical Greenland fjords have lengths $L_f \sim 50$--$100$\,km, cross-sections $W_f \times D_f \sim 5\,\mathrm{km} \times 400\,\mathrm{m}$, giving volumes $V_f \sim 10^{11}$\,m$^3$. Exchange flows are $Q \sim 10^4$\,m$^3$\,s$^{-1}$ \citep{Straneo2012}, so the renewal timescale is
\begin{equation}
	\tau_{\mathrm{fjord}} = \frac{V_f}{Q} \sim \frac{10^{11}}{10^4 \times 3.15 \times 10^7} \approx 0.3\text{--}0.5\ \mathrm{yr}.
	\label{eq:tau_fjord}
\end{equation}
This is fast relative to the dynamic adjustment and is not the rate-limiting step.

\paragraph{Glacier dynamic adjustment.}
\citet{Felikson2017} showed that dynamic thinning propagates inland from the terminus at rates controlled by the glacier's P\'{e}clet number, $\mathrm{Pe} = u H / D_{\mathrm{diff}}$, where $u$ is the depth-averaged velocity, $H$ the ice thickness, and $D_{\mathrm{diff}}$ the effective longitudinal diffusivity set by the membrane stress. The adjustment timescale is $\tau_{\mathrm{dyn}} = \ell_c / c_{\mathrm{pert}}$, where $\ell_c$ is the thinning penetration length and $c_{\mathrm{pert}}$ is the perturbation propagation speed.

\paragraph{Flux-weighted aggregate timescale.}
\citet{Felikson2017} measured thinning penetration lengths for 141 marine-terminating glaciers, finding $\ell_c \sim 2$--$120$\,km with a right-skewed distribution. We approximate this as log-normal with median 15\,km and $\sigma_{\ln \ell} = 0.7$ (90\% range: 5--47\,km). Terminus velocities from MEaSUREs \citep{Mouginot2019} are similarly log-normal with median $\sim$2\,km\,yr$^{-1}$ and $\sigma_{\ln u} = 0.6$. The flux-weighted median is $\tilde{\tau}_{\mathrm{dyn}}^{(D)} = 8.5$\,yr. Adding the fjord renewal time gives a total physical delay estimate of
\begin{equation}
	\delta_{\mathrm{phys}} \approx \tau_{\mathrm{fjord}} + \tilde{\tau}_{\mathrm{dyn}}^{(D)} \approx 0.4 + 8.5 \approx 9\ \mathrm{yr}.
\end{equation}
This is an upper estimate because it assumes the full perturbation must traverse the coupling length before discharge responds; in practice, terminus retreat begins affecting flux on shorter timescales.

\paragraph{Comparison with fitted value.}
Cross-correlation of EN4 subsurface ocean temperature rate versus \citet{Mouginot2019} discharge rate peaks at $\delta = 6$\,yr ($r = 0.84$). BIC on the cumulative demeaned fit selects $\delta = 5$\,yr ($R^2 = 0.995$). Both are shorter than the flux-weighted physical estimate (9\,yr), consistent with the expectation that the aggregate ice-sheet discharge responds before the full dynamic adjustment is complete: terminus retreat and initial speedup occur on shorter timescales than inland thinning propagation.

The range $\delta = 5$--$6$\,yr from cross-correlation and BIC is also consistent with the observed multi-year lag between subsurface ocean temperature peaks and discharge maxima at individual glaciers \citep{Straneo2012, Holland2008}.

%% ================================================================
\section{Ocean--discharge coupling $\gamma$}
\label{sec:gamma}

The parameter $\gamma$ is the rate of sea-level rise per degree of delayed ocean temperature: $\dot{H}_{\mathrm{dyn}} = \gamma\, T_{\mathrm{ocean}}(t - \delta) + r_0$. We derive it from two independent approaches.

\subsection{Observational estimate}
\label{sec:gamma_obs}

Over the calibration window (1972--2018), \citet{Mouginot2019} documented an increase in Greenland discharge of $\Delta D \approx 50$\,Gt\,yr$^{-1}$ above the 1972--1990 baseline. EN4 subsurface ocean temperature (200--500\,m, Greenland peripheral waters) warmed by $\Delta T_{\mathrm{ocean}} \approx 0.4\,^\circ$C over the same period.

In the delay model, the relevant temperature is $T_{\mathrm{ocean}}(2018 - \delta)$. For $\delta = 5$\,yr, the ocean temperature at the effective time (2013) is approximately $\Delta T_{\mathrm{ocean}} \approx 0.35\,^\circ$C. The observational estimate is
\begin{equation}
	\gamma^{\mathrm{obs}} = \frac{\Delta D}{\Delta T_{\mathrm{ocean}}(2018 - \delta)} \approx \frac{50}{0.35} \approx 143\ \mathrm{Gt\,yr^{-1}\,{^\circ C}^{-1}} = 0.39\ \mathrm{mm\,yr^{-1}\,{^\circ C}^{-1}},
\end{equation}
where we use 362\,Gt\,=\,1\,mm GMSL.

\subsection{Bottom-up estimate from glacier dynamics}
\label{sec:gamma_bottomup}

We construct a scaling chain linking ocean temperature to aggregate discharge anomaly through four steps.

\paragraph{Step~1: Submarine melt sensitivity.}
Submarine melting at tidewater termini scales with ocean thermal forcing $T_{\mathrm{ocean}} - T_f$ (where $T_f \approx -2\,^\circ$C is the pressure-dependent freezing point). Plume models and observations give sensitivities of $\partial \dot{m} / \partial T_{\mathrm{ocean}} \sim 5$--$30$\,m\,yr$^{-1}$\,$^\circ$C$^{-1}$, depending on fjord geometry and subglacial discharge \citep{Xu2013, Rignot2016, Jackson2022}. We adopt a central value $\partial \dot{m}/\partial T = 15$\,m\,yr$^{-1}$\,$^\circ$C$^{-1}$.

\paragraph{Step~2: Terminus retreat amplification.}
Submarine melt alone cannot explain the observed retreat rates. Calving amplifies undercutting: melt-induced notch formation, loss of ice m\'{e}lange buttressing, and crevasse propagation produce terminus retreat rates 2--5 times the submarine melt rate anomaly \citep{Benn2007, Luckman2015}. We use an amplification factor $\alpha_c = 3$, giving a terminus retreat rate
\begin{equation}
	\frac{dL}{dT} = \alpha_c \times \frac{\partial \dot{m}}{\partial T} \approx 3 \times 15 = 45\ \mathrm{m\,yr^{-1}\,{^\circ C}^{-1}}.
\end{equation}

\paragraph{Step~3: Velocity sensitivity to retreat.}
\citet{King2020} showed that Greenland's marine-terminating glaciers exhibit a consistent velocity response of 4--5\% speedup per km of terminus retreat, with this relationship holding across the ice sheet and persisting for over a decade. We adopt $\partial \ln u / \partial L = 0.045$\,km$^{-1}$.

\paragraph{Step~4: Integration over the glacier population.}
Total baseline discharge is $D_0 \approx 490$\,Gt\,yr$^{-1}$ \citep{Mouginot2019}. For a sustained $+1\,^\circ$C ocean warming, terminus retreat accumulates at $dL/dT = 45$\,m\,yr$^{-1}$. Over a representative averaging interval $\bar{t}$ (half the calibration window, $\bar{t} \approx 23$\,yr), the mean accumulated retreat is $\bar{L} = (dL/dT) \times \bar{t} \approx 1$\,km, producing a mean fractional speedup of
\begin{equation}
	\frac{\Delta u}{u} = \frac{\partial \ln u}{\partial L} \times \bar{L} \approx 0.045 \times 1 = 4.5\%.
\end{equation}
The corresponding discharge anomaly is
\begin{equation}
	\gamma^{\mathrm{dyn}} = D_0 \times \frac{\Delta u}{u} \approx 490 \times 0.045 \approx 22\ \mathrm{Gt\,yr^{-1}\,{^\circ C}^{-1}} / \bar{D}_{\mathrm{eff}},
\end{equation}
where $\bar{D}_{\mathrm{eff}} \approx 0.17\,^\circ$C for a linear ramp from 0 to 0.35$\,^\circ$C, giving
\begin{equation}
	\gamma^{\mathrm{dyn}} \approx \frac{22}{0.17} \approx 130\ \mathrm{Gt\,yr^{-1}\,{^\circ C}^{-1}} = 0.36\ \mathrm{mm\,yr^{-1}\,{^\circ C}^{-1}}.
	\label{eq:gamma_bottomup}
\end{equation}

\paragraph{Comparison with fitted value.}
The WLS posterior is $\gamma = 0.365$\,mm\,yr$^{-1}$\,$^\circ$C$^{-1}$ (90\% CI: 0.352--0.379). This is in close agreement with both the bottom-up estimate (0.36\,mm\,yr$^{-1}$\,$^\circ$C$^{-1}$) and the observational estimate (0.39\,mm\,yr$^{-1}$\,$^\circ$C$^{-1}$), and well within the uncertainty of the input parameters ($\partial \dot{m}/\partial T$, $\alpha_c$). The two-stage estimation---fixing $\delta$ from cross-correlation before estimating $\gamma$---ensures that the sensitivity is identified by the data rather than traded against the delay.

\subsection{Summary: $\gamma$}

\begin{table}[h]
\centering
\begin{tabular}{lcc}
\toprule
Method & Gt\,yr$^{-1}$\,$^\circ$C$^{-1}$ & mm\,yr$^{-1}$\,$^\circ$C$^{-1}$ \\
\midrule
Observational (delay-corrected) & 143 & 0.39 \\
Bottom-up (glacier dynamics) & 130 & 0.36 \\
WLS posterior (fitted) & 132 & 0.365 \\
\bottomrule
\end{tabular}
\caption{Independent estimates of $\gamma$. All three are mutually consistent within the uncertainties of the input parameters.}
\label{tab:gamma}
\end{table}


%% ================================================================
\section{Secular drift $r_0$}
\label{sec:r0}

The parameter $r_0$ represents a background rate of sea-level rise from Greenland discharge that is independent of recent ocean temperature variability. Two physical sources contribute:

\paragraph{Post-Little Ice Age dynamic adjustment.}
Many of Greenland's outlet glaciers have been retreating from their Little Ice Age (LIA) maxima since the mid-19th century \citep{Kjeldsen2015}. \citet{Kjeldsen2015} estimated that Greenland lost $\sim$1500\,Gt from dynamics between 1900 and 1983, corresponding to $\sim$18\,Gt\,yr$^{-1}$ or $r_0^{\mathrm{LIA}} \approx 0.05$\,mm\,yr$^{-1}$. This ongoing adjustment reflects the slow response of large ice streams to terminus perturbations that originated over a century ago, with the perturbation propagating inland over timescales much longer than $\delta$ \citep{Felikson2017}.

\paragraph{Long-term ocean-forced retreat.}
Pre-satellite ocean temperature reconstructions suggest gradual North Atlantic warming since the early 20th century. The component of discharge driven by this slow warming that is not captured by the delay model (because $T_{\mathrm{ocean}}(t - \delta)$ tracks interannual--decadal variability) would project onto $r_0$. This contribution is difficult to separate from the LIA adjustment and is likely $\sim$5--10\,Gt\,yr$^{-1}$ ($\sim$0.01--0.03\,mm\,yr$^{-1}$).

\paragraph{Comparison with fitted value.}
The combined physical estimate is $r_0 \approx 0.06$--$0.08$\,mm\,yr$^{-1}$. The WLS posterior is $r_0 = 0.104$\,mm\,yr$^{-1}$ (90\% CI: 0.099--0.108), approximately 1.4 times the physical estimate. The modest excess likely reflects additional secular processes not captured by the two physical sources above, such as pre-satellite ocean warming trends that are not well represented in the EN4 reconstruction before Argo coverage ($\sim$2004).

%% ================================================================
\section{Synthesis}
\label{sec:synthesis}

All three parameters can be independently estimated from glacier dynamics and observations without reference to the WLS calibration:

\begin{table}[h]
\centering
\begin{tabular}{lccc}
\toprule
Parameter & Physical estimate & WLS posterior & Ratio \\
\midrule
$\delta$ (yr) & 5--9 & 5 (BIC) & 1.0 \\
$\gamma$ (mm\,yr$^{-1}$\,$^\circ$C$^{-1}$) & 0.36--0.39 & 0.365\,[0.352, 0.379] & 1.0 \\
$r_0$ (mm\,yr$^{-1}$) & 0.06--0.08 & 0.104\,[0.099, 0.108] & 1.4 \\
\bottomrule
\end{tabular}
\caption{Comparison of physically derived and fitted parameter values. The ratio column gives (fitted)/(physical midpoint).}
\label{tab:synthesis}
\end{table}

\noindent The agreement for $\delta$ and $\gamma$ is within the uncertainty of the physical estimates. The factor-of-1.4 excess in $r_0$ is modest and likely reflects secular processes not fully represented in the two physical sources above.

%% ================================================================
\section{Sensitivity to RCM choice for SMB}
\label{sec:rcm_sensitivity}

Our Greenland SMB observations are taken from \citet{Mouginot2019}, which uses RACMO2.3p2 \citep{Noel2018} as its sole RCM for surface mass balance. To assess sensitivity to this choice, we compare against the \citet{Mankoff2021} product, which averages SMB from three RCMs---HIRHAM/HARMONIE \citep{Langen2017}, MAR v3.12 \citep{Fettweis2020}, and RACMO---over the overlapping period (1972--2018).

The two cumulative SMB time series agree within observational uncertainties throughout the satellite era and the earlier Mouginot record. The Mankoff three-RCM average is systematically slightly lower (less mass loss) than Mouginot's RACMO-only SMB, consistent with the GrSMBMIP intercomparison \citep{Fettweis2020}, which found MAR produces $\sim$15\,Gt\,yr$^{-1}$ higher mean SMB than RACMO over 1980--2012 (372 vs.\ 357\,Gt\,yr$^{-1}$, a $\sim$4\% difference). The Mankoff values also track more closely to the RCM-independent GRACE-minus-discharge implied SMB, suggesting the multi-RCM average may marginally reduce individual model biases.

Because the RACMO-only versus three-RCM difference ($\sim$15\,Gt\,yr$^{-1}$, or $\sim$0.04\,mm\,yr$^{-1}$ SLE) is small relative to both observational uncertainties and our other structural uncertainties (discharge model, ocean temperature reconstruction), the RCM choice does not materially affect our hindcast or projections. We retain RACMO via \citet{Mouginot2019} for consistency with the discharge calibration, which also uses the Mouginot dataset.

\bibliographystyle{agufull08}
\bibliography{references}

\end{document}
